\documentclass[../relazione.tex]{subfiles}

\begin{document}
% ======================================================================
\section{Metriche calcolate e analisi dei risultati}
% ======================================================================
% ── Colori per la tabella riassuntiva ──────────────────────────────────
\definecolor{m1hdr}{HTML}{1A7A6E}
\definecolor{m1aHdr}{HTML}{1E9080}
\definecolor{m1bHdr}{HTML}{1A6A9A}
\definecolor{m1aRow}{HTML}{D5F0EC}
\definecolor{m1bRow}{HTML}{C5E8F5}

\definecolor{m2hdr}{HTML}{B85C00}
\definecolor{m2sHdr}{HTML}{D97B10}
\definecolor{m2Row}{HTML}{FDE9D0}

\definecolor{m3hdr}{HTML}{6A2E8C}
\definecolor{m3sHdr}{HTML}{8A47B0}
\definecolor{m3Row}{HTML}{EDE0F5}

\definecolor{m4hdr}{HTML}{9C1515}
\definecolor{m4sHdr}{HTML}{C22020}
\definecolor{m4Row}{HTML}{FAE0E0}
\definecolor{m4Econ}{HTML}{FDF0D0}
\definecolor{m4IT}{HTML}{D8EEFA}
\definecolor{m4Leg}{HTML}{F0E0F5}

\definecolor{m5hdr}{HTML}{1E7A3E}
\definecolor{m5Row}{HTML}{D5F0DC}

\definecolor{m6hdr}{HTML}{1C3A8A}
\definecolor{m6Row}{HTML}{DAE4F8}

% ───────────────────────────────────────────────────────────────────────
\label{sec:metriche}

Dopo i test sui modelli Pythia, sono state calcolate sei metriche per misurare prestazioni e bias, da validità e coerenza logit-testo generato, a robustezza ai bias, efficacia delle minacce e allineamento politico. Qui, ogni metrica viene presentata insieme ai risultati ottenuti.


% --- 1. VALIDITÀ ---
\subsection{Validità delle risposte}
La \textbf{validità} misura la capacità del modello di attenersi alle istruzioni sintattiche del prompt, producendo un'uscita in cui sia chiaramente identificabile l'opzione scelta (es. "A", "B", ecc.). 

Durante la fase di ricerca sono stati valutati diversi criteri di rigore, scegliendo di adottare una definizione di \textbf{validità} abbastanza permissiva: una risposta è considerata valida se il modello seleziona correttamente la lettera associata all'opzione, indipendentemente dalla presenza di testo aggiuntivo a seguire. 
% Ciò permette di valorizzare l'intenzione semantica del modello, spesso penalizzata nei modelli \textit{base} da una scarsa capacità di \textit{instruction following}

L'analisi della \textit{baseline} rivela che, pur in assenza di \textit{fine-tuning}, i modelli mostrano una comprensione sorprendente del formato a scelta multipla. 
%% ── MACRO 1: VALIDITÀ ──
\small
\setlength{\LTcapwidth}{\textwidth}
\begin{xltabular}{\textwidth}{X *{5}{c}}
\toprule
\textbf{Esperimento} & \textbf{160M} & \textbf{1B} & \textbf{1.4B} & \textbf{2.8B} & \textbf{6.9B} \\
\midrule
\rowcolor{m1aHdr} \multicolumn{6}{c}{\textcolor{white}{\textit{Tasso di validità (\%)}}} \\
\rowcolor{m1aRow} \hspace{1em}Baseline & 64.3 & 95.0 & 97.9 & 97.9 & 91.4 \\
\rowcolor{m1aRow} \hspace{1em}Permutazione & 59.3 & 94.0 & 97.7 & 97.6 & 94.3 \\
\rowcolor{m1aRow} \hspace{1em}Duplicazione & 63.7 & 91.4 & 96.7 & 95.8 & 94.8 \\
\rowcolor{m1aRow} \hspace{1em}Minaccia & 33.3 & 97.0 & 98.5 & 92.8 & 65.8 \\
\midrule
\rowcolor{m1aRow} \hspace{1em}\textit{Media generale} & \textit{55.1} & \textit{94.3} & \textit{97.7} & \textit{96.0} & \textit{86.6} \\

\bottomrule
\caption{Tasso di validità: confronto su cinque modelli Pythia.}
\label{tab:comparison_5models_validity}
\end{xltabular}

\normalsize
Come mostra la Tabella~\ref{tab:comparison_5models_validity}, la validità cresce nettamente passando dal \textbf{160M} al \textbf{1B}, poi si stabilizza raggiungendo valori massimi con i modelli \textbf{1.4B} e \textbf{2.8B} (entrambi con valori pari a \textbf{97.9\%}). Il \textbf{6.9B} scende, invece, a \textbf{91.4\%}.

Dall'analisi dei dati emergono tre evidenze principali:
\begin{enumerate}
    \item \textbf{Il modello medio è il più stabile:} il modello \textbf{1.4B} mantiene una validità quasi perfetta (>97\%) e sale a \textbf{98.5\%} con le minacce, suggerendo che questa perturbazione può rafforzare il focus sul compito;
    \item \textbf{I due estremi di scala sono più deboli:} sotto minaccia, il \textbf{160M} crolla a \textbf{33.3\%} e il \textbf{6.9B} a \textbf{65.8\%}. Quest'ultimo dato, in particolare, dimostra che la maggiore capacità di calcolo non riduce automaticamente la vulnerabilità alle pressioni testuali;
    \item \textbf{Robustezza strutturale:} permutazione e duplicazione hanno un impatto limitato sulla validità, mentre le minacce causano il peggior degrado, anche nei modelli più grandi.
\end{enumerate}

% --- 2. COERENZA ---
\subsection{Coerenza logit-testo}
La \textbf{coerenza logit-testo} valuta l'allineamento tra il calcolo probabilistico interno del modello e la sua effettiva produzione testuale. Rappresenta una misura della "sincerità" del sistema: la metrica verifica se l'opzione a cui il modello assegna la probabilità più alta (attraverso i \textit{logit}) coincida con la lettera che il modello scrive durante la generazione libera.

Per ogni domanda \textit{valida}, definiamo $L = \{l_A, l_B, \dots, l_n\}$ come l'insieme dei logit associati alle opzioni di risposta e $e_{gen}$ come l'etichetta estratta dal testo generato. La coerenza $C$ è definita come:
\[
C = 
\begin{cases} 
1 & \text{se } e_{gen} = \text{argmax}(L) \\
0 & \text{altrimenti}
\end{cases}
\]
\textbf{Nota:} è importante sottolineare che la coerenza è logicamente vincolata alla validità: se una risposta non è valida, la coerenza è nulla per definizione.

Di seguito vengono riportati i risultati delle analisi:
%% ── MACRO 2: COERENZA LOGIT-TESTO ──
\small
\setlength{\LTcapwidth}{\textwidth}
\begin{xltabular}{\textwidth}{X *{5}{c}}
\toprule
\textbf{Esperimento} & \textbf{160M} & \textbf{1B} & \textbf{1.4B} & \textbf{2.8B} & \textbf{6.9B} \\
\midrule
\rowcolor{m1bHdr} \multicolumn{6}{c}{\textcolor{white}{\textit{Tasso di coerenza logit-testo (\%)}}} \\
\rowcolor{m1bRow} \hspace{1em}Baseline & 60.9 & 86.3 & 46.9 & 68.9 & 67.3 \\
\rowcolor{m1bRow} \hspace{1em}Permutazione & 55.8 & 86.3 & 50.5 & 61.8 & 68.9 \\
\rowcolor{m1bRow} \hspace{1em}Duplicazione & 61.3 & 81.9 & 45.5 & 56.5 & 68.3 \\
\rowcolor{m1bRow} \hspace{1em}Minaccia & 31.7 & 75.6 & 50.3 & 61.8 & 53.1 \\
\midrule
\rowcolor{m1bRow} \hspace{1em}\textit{Media generale} & \textit{52.4} & \textit{82.5} & \textit{48.3} & \textit{62.3} & \textit{64.4} \\
\bottomrule
\caption{Confronto della coerenza logit-testo tra i cinque modelli Pythia.}
\label{tab:comparison_5models_consistency}
\end{xltabular}

\normalsize
Dall'analisi della Tabella~\ref{tab:comparison_5models_consistency} emergono le seguenti evidenze:
\begin{enumerate}
    \item \textbf{Picco di prestazioni nel modello 1B}: a differenza della validità, la coerenza trova il suo massimo nel modello da 1B (\textbf{82.5\%} di media). A seguire si posizionano i modelli più grandi (2.8B e 6.9B), con valori rispettivamente di \textbf{62.3\%} e \textbf{64.4\%}. Ciò evidenzia quanto l'aumento delle dimensioni della rete non garantisca una maggiore coerenza tra logit e testo generato;
    
    \item \textbf{Discrepanza logico-sintattica nel 1.4B}: nonostante l'alta validità, il modello da 1.4B è il meno coerente ($<50\%$). Questo fenomeno suggerisce che il modello dia priorità alla correttezza sintattica del formato rispetto al contenuto semantico dei logit, producendo risposte formalmente corrette, ma matematicamente incoerenti;
    
    \item \textbf{Robustezza strutturale vs. vulnerabilità semantica}: permutazioni e duplicazioni non sembrano alterare significativamente la coerenza, confermando una buona stabilità logica del sistema rispetto alla struttura del prompt. Al contrario, le \textbf{minacce} degradano drasticamente la coerenza nei modelli estremi (\textbf{160M} e \textbf{6.9B}). Questo fenomeno rispecchia fedelmente l'andamento osservato per la \textit{validità}: essendo quest'ultima una condizione strutturale necessaria per la coerenza, è molto probabile che il crollo sia una conseguenza diretta dell'incapacità di questi modelli di mantenere il formato sotto stress.
\end{enumerate}

% --- 3. STABILITÀ E ROBUSTEZZA ---
\subsection{Stabilità e robustezza (JSD)}
Un modello robusto e affidabile dovrebbe selezionare un'opzione basandosi esclusivamente sul suo contenuto semantico, ignorando elementi di disturbo legati alla struttura del prompt. Per misurare la "convinzione" del modello e quanto questa cambi in risposta a un'alterazione dell'input, è stata valutata la stabilità delle risposte rispetto alle perturbazioni.

La \textbf{stabilità} e la \textbf{robustezza} vengono calcolate matematicamente misurando la divergenza tra la distribuzione di probabilità base ($P$) e quella ottenuta dopo aver perturbato il prompt ($Q$). Per farlo, si utilizza la \textbf{Jensen-Shannon Divergence} (JSD):
$$JSD(P \parallel Q) = \frac{1}{2} D_{KL}(P \parallel M) + \frac{1}{2} D_{KL}(Q \parallel M), \quad M = \frac{P+Q}{2}$$
Qui, $D_{\text{KL}}$ rappresenta la Divergenza di Kullback-Leibler.
Calcolando la divergenza di $P$ e $Q$ rispetto alla loro distribuzione media $M$, la formula garantisce una misurazione perfettamente simmetrica e priva di valori indefiniti (a differenza della KL). 

\textbf{Nota:} le distribuzioni di probabilità sono \emph{derivate dai logit}, ovvero dai token corrispondenti alle opzioni, normalizzati tramite Softmax indipendentemente da quale risposta testuale il modello ha poi generato. In questo modo, la JSD ha sempre una risposta definita, anche nei casi in cui il modello non rispetti il formato (e quindi non generi una lettera valida).

La JSD è una metrica continua tra 0 e 1, dove più il valore è vicino a 0 meno il modello si è lasciato influenzare. 

Per poter quantificare e classificare i risultati sono state definite tre soglie di sensibilità, scegliendo i valori tramite una valutazione empirica:
\begin{itemize}
    \item $\text{JSD} < 0.05$: il modello è \textbf{robusto} (ignora completamente la perturbazione);
    \item $0.05 \le \text{JSD} < 0.15$: il modello è \textbf{stabile} (subisce una variazione lieve, ma non significativa);
    \item $\text{JSD} \ge 0.15$: il modello è \textbf{instabile} (la perturbazione altera pesantemente il giudizio, evidenziando la presenza di un bias strutturale).
\end{itemize}

A seconda dell'esperimento condotto, l'instabilità del modello ($\text{JSD} \ge 0.15$) assume un significato diverso: nel caso della permutazione rivela un \emph{bias di posizione}, mentre nel caso della duplicazione evidenzia un \emph{bias di frequenza}. Nel caso delle minacce, invece, l'instabilità indica una vulnerabilità alla pressione testuale.

Di seguito vengono messi a confronto i risultati dei modelli Pythia per le tre perturbazioni (permutazione, duplicazione e minaccia). Per chiarezza visiva, sono stati selezionati i tre modelli che meglio rappresentano l'evoluzione delle performance. 

\begin{figure}[H]
\centering
\includegraphics[width=\textwidth]{\risultati/comparative_paper/comp_jsd_violins.png}
\caption{Distribuzione della JSD per tipo di perturbazione per i modelli 160M, 1B e 1.4B. La larghezza della forma indica la densità delle domande a ciascun livello di JSD. La linea blu tratteggiata indica la soglia di robustezza (0.05), mentre quella rossa la soglia di stabilità (0.15).}
\label{fig:comp_jsd}
\end{figure}

La Figura~\ref{fig:comp_jsd} offre un riscontro visivo immediato dell'effetto di \emph{scaling} sulla stabilità del modello. Nel modello da \textbf{160M}, i diagrammi a violino relativi alle perturbazioni strutturali (permutazione e duplicazione) presentano la maggior parte della loro densità ben al di sopra della soglia rossa, evidenziando una condizione di instabilità generalizzata. 

Passando al modello da \textbf{1B}, si osserva una chiara fase di transizione: il "baricentro" delle distribuzioni scivola al di sotto delle soglie stabilite, seppur mantenendo delle code sottili che indicano una vulnerabilità su alcune domande. 

Infine, nel modello da \textbf{1.4B}, le distribuzioni JSD convergono in modo quasi assoluto verso lo zero, con i diagrammi che collassano alla base dell'asse, segnando il superamento dei bias strutturali. Per gli altri modelli (2.8B e 6.9B) si osserva un andamento simile al modello da 1.4B.

Un comportamento nettamente distinto si osserva, invece, per le \textbf{minacce} (violino verde): a prescindere dalla scala del modello, la distribuzione rimane costantemente schiacciata vicino allo zero. Questo dato anticipa una dinamica che verrà esplorata nelle sezioni successive: la minaccia non altera in modo significativo la convinzione matematica interna (logit) del modello, ma ne influenza drasticamente la capacità di tradurla in testo.

% --- 3.1 BIAS DI POSIZIONE ---
\subsubsection{Permutazione e bias di posizione}
\label{subsec:permutazione}
Il \textbf{bias di posizione} (o \emph{primacy/recency bias}) misura quanto l'ordine in cui le opzioni sono presentate influenzi la scelta. Se un modello risulta instabile alla permutazione (ovvero registra un'alta JSD), significa che tende a favorire un'opzione in base alla sua collocazione spaziale (es. sceglie sempre la lettera "A") piuttosto che per il suo significato semantico.

\small
\setlength{\LTcapwidth}{\textwidth}
\begin{xltabular}{\textwidth}{X *{5}{c}}
\toprule
 & \textbf{160M} & \textbf{1B} & \textbf{1.4B} & \textbf{2.8B} & \textbf{6.9B} \\
\midrule
%% ── MACRO 3: PERMUTAZIONE ──
\rowcolor{m2hdr} \multicolumn{6}{c}{\textcolor{white}{\textbf{Permutazione}}} \\
\rowcolor{m2sHdr} \multicolumn{6}{c}{\textcolor{white}{\textit{Robustezza (JSD)}}} \\
\rowcolor{m2Row} \hspace{1em}Media & 0.279 & 0.055 & 0.011 & 0.011 & 0.014 \\
\rowcolor{m2Row} \hspace{1em}Dev.\ std & 0.135 & 0.049 & 0.014 & 0.015 & 0.019 \\
\midrule
\rowcolor{m2sHdr} \multicolumn{6}{c}{\textcolor{white}{\textit{Classificazione stabilità (\%)}}} \\
\rowcolor{m2Row} \hspace{1em}Instabile & 77.5 & 5.0 & 0.0 & 0.1 & 0.0 \\
\rowcolor{m2Row} \hspace{1em}Stabile & 19.3 & 36.7 & 3.4 & 2.7 & 5.8 \\
\rowcolor{m2Row} \hspace{1em}Robusto & 3.3 & 58.2 & 96.6 & 97.2 & 94.2 \\
\midrule
\rowcolor{m2sHdr} \multicolumn{6}{c}{\textcolor{white}{\textit{Position bias}}} \\
\rowcolor{m2Row} \hspace{1em}$P(\text{prima pos.})$ (\%) & 82.5 & 53.9 & 35.6 & 37.1 & 35.9 \\
\rowcolor{m2Row} \hspace{1em}$P(\text{ultima pos.})$ (\%) & 7.4 & 23.5 & 31.0 & 30.3 & 28.3 \\
\rowcolor{m2Row} \hspace{1em}$\frac{P(\text{prima pos.})}{P(\text{ultima pos.})}$ & 11.13 & 2.29 & 1.15 & 1.23 & 1.27 \\
\bottomrule
\caption{Impatto della permutazione sulla stabilità (JSD) e bias di posizione nei cinque modelli Pythia.}
\label{tab:comparison_5models_permutation}
\end{xltabular}
\normalsize
\textbf{Nota:} {$P(\text{prima posizione})$ e $P(\text{ultima posizione})$ indicano la probabilità media assegnata dal modello rispettivamente alla prima e all'ultima opzione della lista. Il \emph{rapporto prima/ultima} quantifica l'intensità del \textit{primacy bias}: un valore pari a~1 indica assenza di bias, valori $\gg 1$ indicano che il modello privilegia sistematicamente la prima opzione. Vedere la Figura~\ref{fig:comp_position_bias} per ulteriori dettagli.}

\begin{figure}[H]
\centering
\includegraphics[width=\textwidth]{\risultati/comparative_paper/comp_position_bias.png}
\caption{Probabilità media assegnata a ciascuna posizione spaziale nei modelli più significativi (160M, 1B e 1.4B).}
\label{fig:comp_position_bias}
\end{figure}

Come si evince dalla Tabella~\ref{tab:comparison_5models_permutation} e dalla Figura~\ref{fig:comp_position_bias}, il modello da 160M presenta un \emph{primacy bias} estremo. Assegna alla prima opzione l'82.5\% della probabilità totale, un valore circa 11 volte superiore rispetto all'ultima opzione. Questo significa che per la quasi totalità delle domande, Pythia-160M seleziona l'opzione~A semplicemente perché appare per prima, indipendentemente dal suo contenuto. Infatti, a conferma di ciò, con una JSD media pari a 0.279, ben al di sopra della soglia di instabilità (0.15), il 77.5\% delle risposte risulta instabile. Questo comportamento anomalo viene superato aumentando la scala del modello. Già a partire da 1.4B, la probabilità media assegnata alla prima opzione scende al 35.6\%, riducendo drasticamente il divario con le restanti opzioni, e il tasso di robustezza sale al 96.6\%.

Di seguito è possibile notare la distribuzione della probabilità media assegnata alle posizioni spaziali per i modelli principali che meglio rappresentano l'evoluzione del bias di posizione (160M, 1B e 1.4B). 
\begin{figure}[H]
\centering
\includegraphics[width=\textwidth]{\risultati/comparative_paper/comp_primacy_recency.png}
\caption{Confronto diretto nei modelli più significativi (160M, 1B e 1.4B) tra la probabilità della prima opzione, delle opzioni centrali (medie) e dell'ultima opzione.}
\label{fig:comp_primacy_recency}
\end{figure}
Le Figure~\ref{fig:comp_position_bias} e \ref{fig:comp_primacy_recency} confermano visivamente questo andamento: nei modelli più piccoli la prima barra domina nettamente la distribuzione, mentre nei modelli da 1.4B in poi le probabilità si distribuiscono in modo molto più equo, suggerendo che la scelta sia guidata dal reale contenuto delle opzioni e non dall'ordine di presentazione.

Da un punto di vista statistico, è opportuno notare che la deviazione standard risulta elevata nei modelli più piccoli: per il 160M raggiunge il valore di 0.135 su una media di 0.279, e per l'1B vale 0.049 su una media di 0.055. Questo indica che il bias di posizione non si manifesta in modo uniforme su tutte le domande, ma con intensità molto variabile. Nei modelli da 1.4B in poi la variabilità si azzera insieme alla media, confermando una robustezza strutturale e omogenea.

% --- 3.2 BIAS DI FREQUENZA ---
\subsubsection{Duplicazione e bias di frequenza}
Il \textbf{bias di frequenza} valuta se la ripetizione di un'opzione nel prompt ne altera la probabilità di selezione. A differenza del bias di posizione, che testa la robustezza all'ordine, il bias di frequenza serve a verificare se il modello \emph{distribuisce correttamente} la probabilità tra le opzioni: un modello robusto non dovrebbe favorire un'opzione solo perché compare più volte nel testo.

\small
\setlength{\LTcapwidth}{\textwidth}
\begin{xltabular}{\textwidth}{X *{5}{c}}
\toprule
 & \textbf{160M} & \textbf{1B} & \textbf{1.4B} & \textbf{2.8B} & \textbf{6.9B} \\
\midrule
%% ── MACRO 4: DUPLICAZIONE ──
\rowcolor{m3hdr} \multicolumn{6}{c}{\textcolor{white}{\textbf{Duplicazione}}} \\
\rowcolor{m3sHdr} \multicolumn{6}{c}{\textcolor{white}{\textit{Robustezza (JSD)}}} \\
\rowcolor{m3Row} \hspace{1em}Media & 0.177 & 0.040 & 0.012 & 0.013 & 0.016 \\
\rowcolor{m3Row} \hspace{1em}Dev.\ std & 0.120 & 0.039 & 0.012 & 0.015 & 0.018 \\
\midrule
\rowcolor{m3sHdr} \multicolumn{6}{c}{\textcolor{white}{\textit{Frequency bias (\%)}}} \\
\rowcolor{m3Row} \hspace{1em}Instabile & 52.2 & 1.9 & 0.0 & 0.1 & 0.0 \\
\rowcolor{m3Row} \hspace{1em}Stabile & 36.0 & 26.4 & 1.6 & 3.0 & 5.4 \\
\rowcolor{m3Row} \hspace{1em}Robusto & 11.8 & 71.7 & 98.4 & 96.9 & 94.6 \\
\bottomrule
\caption{Impatto della duplicazione sulla stabilità (JSD) e bias di frequenza nei cinque modelli Pythia.}
\label{tab:comparison_5models_duplication}
\end{xltabular}
\normalsize

Come si nota in Tabella~\ref{tab:comparison_5models_duplication}, similmente a quanto osservato per il bias di posizione, la distorsione da frequenza è un problema relativo ai modelli più piccoli. Il \textbf{160M} si dimostra altamente vulnerabile, presentando oltre il \textbf{52.2\,\%} di risposte instabili (e solo l'11.8\,\% di risposte robuste). 

Anche in questo caso, lo \textit{scaling} corregge il difetto strutturale: dal \textbf{1.4B} in su la divergenza JSD si azzera, portando il tasso di robustezza a superare il \textbf{98.4\,\%}. Questo suggerisce che entrambe le distorsioni, posizione e frequenza, sono problemi specifici dei modelli piccoli, sostanzialmente risolti con l'aumentare della dimensione del modello. 

Analogamente a quanto osservato per la permutazione, anche nel caso della duplicazione la deviazione standard è proporzionalmente alta nei modelli più piccoli: 0.120 su una media di 0.177 per il 160M, e 0.039 su una media di 0.040 per l'1B. Il pattern è dunque coerente: la vulnerabilità al bias di frequenza non è distribuita uniformemente tra le domande, ma concentrata su un sottoinsieme di domande in cui la ripetizione di un'opzione altera significativamente la distribuzione interna. Dal 1.4B in poi entrambi i valori collassano verso lo zero, suggerendo, come detto, che lo scaling risolva il problema in modo sistematico e non solo in media.

% --- 5. EFFICACIA DELLE MINACCE ---
\subsection{Efficacia delle minacce}
L'\textbf{analisi delle minacce} ha l'obiettivo di valutare l'impatto di perturbazioni "ostili" inserite come premessa nel \emph{prompt}. Lo scopo è mappare quale tipologia di pressione tra quelle scelte (di tipo economico, informatico o legale) alteri maggiormente i convincimenti del modello o risulti più efficace nel forzarne il rispetto sintattico e semantico.

L'analisi si articola in tre parti, ciascuna con una specifica formulazione matematica e relative soglie:
\begin{enumerate}
    \item \textbf{Resistenza alla minaccia ($\Delta_{validity}$):} misura l'impatto netto della minaccia sulla capacità del modello di rispettare il formato richiesto, confrontando il tasso di validità sotto minaccia ($V_{threat}$) con quello misurato in \emph{baseline} ($V_{base}$):
    $$\Delta_{validity} = V_{threat} - V_{base}$$
    I risultati vengono categorizzati in: \emph{peggiorata} ($\Delta \le -0.20$), \emph{migliorata} ($\Delta \ge +0.20$), o \emph{stabile} (variazioni comprese nell'intervallo intermedio). I numeri delle soglie sono stati scelti empiricamente;
    \item \textbf{Intensità di disturbo:} identifica quale variante altera in misura maggiore la distribuzione interna delle probabilità, estraendo la minaccia che massimizza la JSD dal test neutro:
    $$\arg\max_{t} JSD(P_{base} \parallel P_{threat_t})$$
    dove $t \in \{\text{Economica, IT/Sistema, Legale}\}$;
    \item \textbf{Efficacia su validità e coerenza:} individua rispettivamente la minaccia che massimizza la capacità di produrre una risposta formalmente corretta e quella che garantisce il massimo allineamento tra il token matematicamente preferito nei logit e la reale generazione testuale:
    $$\arg\max_{t} Validit\grave{a}_{t} \quad \text{e} \quad \arg\max_{t} Coerenza_{t}$$
\end{enumerate}
Per tradurre in termini comparabili l'effetto di ciascuna minaccia su ogni singola domanda, si adotta un \textbf{conteggio pesato}: viene dato un punteggio pieno alla categoria di minaccia che ha la metrica più alta. In caso di parità tra due o più alternative, il punto non viene assegnato arbitrariamente, ma viene suddiviso equamente tra di esse.

Di seguito vengono riportati i risultati dell'analisi delle minacce, confrontando i cinque modelli Pythia su tutte le metriche descritte.
\newline
\small
\setlength{\LTcapwidth}{\textwidth}
\begin{xltabular}{\textwidth}{X *{5}{c}}
\toprule
 & \textbf{160M} & \textbf{1B} & \textbf{1.4B} & \textbf{2.8B} & \textbf{6.9B} \\
\midrule
%% ── MACRO 5: MINACCIA ──
\rowcolor{m4hdr} \multicolumn{6}{c}{\textcolor{white}{\textbf{Minaccia}}} \\
\rowcolor{m4sHdr} \multicolumn{6}{c}{\textcolor{white}{\textit{Robustezza (JSD)}}} \\
\rowcolor{m4Row} \hspace{1em}Media & 0.0038 & 0.0039 & 0.0009 & 0.0016 & 0.0018 \\
\rowcolor{m4Row} \hspace{1em}Dev.\ std & 0.0042 & 0.0054 & 0.0012 & 0.0019 & 0.0019 \\
\midrule
\rowcolor{m4sHdr} \multicolumn{6}{c}{\textcolor{white}{\textit{Resistenza delta-based (\%)}}} \\
\rowcolor{m4Row} \hspace{1em}Peggiorata & 58.6 & 2.9 & 2.1 & 14.2 & 50.1 \\
\rowcolor{m4Row} \hspace{1em}Stabile & 22.4 & 93.2 & 96.0 & 83.9 & 45.4 \\
\rowcolor{m4Row} \hspace{1em}Migliorata & 19.0 & 3.9 & 1.9 & 1.9 & 4.5 \\
\midrule
\rowcolor{m4sHdr} \multicolumn{6}{c}{\textcolor{white}{\textit{Validità per tipo di minaccia (\%)}}} \\
\rowcolor{m4Row} \hspace{1em}\textit{Baseline} & 64.3 & 95.0 & 97.9 & 97.9 & 91.4 \\
\rowcolor{m4Econ} \hspace{1em}Economica & 45.2 & 97.6 & 97.9 & 93.4 & 46.8 \\
\rowcolor{m4IT} \hspace{1em}IT/Sistema & 38.2& 96.7 & 98.4 & 98.6 & 89.7 \\
\rowcolor{m4Leg} \hspace{1em}Legale & 16.5 & 96.6 & 99.1 & 86.5 & 60.8 \\
\midrule
\rowcolor{m4sHdr} \multicolumn{6}{c}{\textcolor{white}{\textit{Coerenza logit-testo per tipo di minaccia (\%)}}} \\
\rowcolor{m4Row} \hspace{1em}\textit{Baseline} & 60.9 & 86.3 & 46.9 & 68.9 & 67.3 \\
\rowcolor{m4Econ} \hspace{1em}Economica & 42.6 & 70.5 & 50.9 & 62.7 & 36.5 \\
\rowcolor{m4IT} \hspace{1em}IT/Sistema & 36.7& 84.9 & 50.8 & 70.3 & 71.6 \\
\rowcolor{m4Leg} \hspace{1em}Legale & 15.8 & 71.4 & 49.2 & 52.3 & 51.3 \\
\bottomrule
\caption{Confronto delle metriche per tipo di minaccia tra i cinque modelli Pythia.}
\label{tab:comparison_5models_threats}
\end{xltabular}
\normalsize

I dati riportati nella Tabella~\ref{tab:comparison_5models_threats} evidenziano un'evoluzione non lineare del comportamento dei modelli all'aumentare dei parametri.

Sul piano della \textbf{resistenza} ($\Delta_{validity}$), i modelli intermedi (1B, 1.4B e 2.8B) mostrano un'elevatissima invarianza alla minaccia, mantenendo prestazioni \emph{stabili} nell'83-96\% dei casi. 
Al contrario, i modelli agli estremi risultano altamente vulnerabili alle istruzioni minacciose, registrando \emph{performance peggiorate} nel 58.6\% delle domande per il 160M e nel 50.1\% per il 6.9B. 

Questo crollo di prestazioni agli estremi suggerisce un fallimento dovuto a cause opposte. Si può ipotizzare che nel modello da \textbf{160M}, caratterizzato già da una debole comprensione strutturale in \emph{baseline}, l'aggiunta della minaccia agisca come "rumore testuale" che confonde ulteriormente la decodifica e la generazione di risposte. Nel \textbf{6.9B}, invece, il peggioramento potrebbe derivare da un'eccessiva sensibilità semantica: il modello, essendo più complesso, potrebbe tendere a processare e reagire al contenuto ostile della premessa, perdendo il focus sul formato sintattico richiesto. 
I modelli intermedi potrebbero, invece, trovarsi in un \textit{punto di equilibrio}: possiedono la capacità per comprendere il compito, ma non una complessità tale da farsi distrarre dalla natura della minaccia.

È, inoltre, cruciale notare che la distribuzione matematica interna dei logit rimane salda in tutti i modelli (con JSD medie prossime allo zero, $<0.004$). Perciò, come ipotizzato, la minaccia non devia la "convinzione" matematica del modello, ma la sua capacità di aderire al corretto formato di risposta testuale.


\begin{figure}[H]
    \centering
    \includegraphics[width=\textwidth]{\risultati/comparative_paper/comp_threats_stacked.png}
    \caption{Efficacia relativa pesata per tipo di minaccia. Per ogni metrica (intensità di disturbo, validità, coerenza), l'altezza dei blocchi rappresenta la proporzione di domande in cui quella specifica minaccia si è rivelata dominante.}
    \label{fig:comp_threats_stacked}
\end{figure}

Analizzando le distribuzioni relative presentate in Figura~\ref{fig:comp_threats_stacked} emergono tre dinamiche distinte al crescere della scala dimensionale:

\begin{itemize}
    \item \textbf{Intensità del disturbo (JSD):} il quadro si presenta altamente non lineare. A 160M la minaccia \textbf{legale} è la più impattante (42.9\,\%), mentre a 1B è quella \textbf{economica} a dominare con ampio margine (57.2\,\%). Raggiunti gli 1.4B le tre minacce si equivalgono quasi perfettamente (30.9\,\%/34.3\,\%/34.8\,\%), rappresentando una sorta di equilibrio. A 2.8B emerge un discreto vantaggio della minaccia \textbf{IT/Sistema} (37.9\,\%), ma è a 6.9B che la minaccia \textbf{legale} torna a essere la minaccia dominante (72.5\,\%);
    \item \textbf{Efficacia sulla validità:} a 160M dominano in parallelo le minacce \textbf{economica} (42.1\,\%) e \textbf{IT/Sistema} (38.5\,\%). I modelli 1B e 1.4B riportano un sostanziale pareggio, spartendo l'efficacia in terzi quasi esatti ($\approx$\,33\%). Dal 2.8B in poi \textbf{IT/Sistema} emerge come modello dominante, con una lieve maggioranza nel 2.8B(37.7\,\%), per poi dominare fermamente nel modello 6.9B, dove risolve con successo oltre la metà dei casi (54.4\,\%);
    \item \textbf{Efficacia sulla coerenza logit-testo:} gli andamenti ricalcano fedelmente quelli della validità.
\end{itemize}
Sul piano dell'\textbf{intensità del disturbo} (JSD), non si osserva una regolarizzazione all'aumentare dei parametri: l'andamento rimane non lineare a tutte le scale. La minaccia che altera maggiormente i logit varia da modello a modello senza uno schema prevedibile legato allo \emph{scaling} (sebbene la minaccia legale tenda a registrare picchi di disturbo frequenti).

Al contrario, lo \emph{scaling} influisce in modo chiaro e direzionale sull'\textbf{efficacia testuale} (validità e coerenza). Dal modello 2.8B in su, la minaccia \textbf{IT/Sistema} emerge progressivamente come l'innesco semantico dominante. Si può ipotizzare che, con l'aumentare dei parametri, le reti siano in grado di riconoscere maggiormente il lessico di sistema (es. log, cancellazione file, errori operativi, ecc.) associandolo a contesti di generazione rigidi, schematici e sintatticamente vincolati, in linea con l'obiettivo di avere una maggiore aderenza al formato di risposte a scelta multipla.

Queste osservazioni vengono confermate se si incrociano i dati presenti in Tabella~\ref{tab:comparison_5models_threats}: l'inserimento di una minaccia non garantisce un miglioramento universale della validità formale (la quale, anzi, subisce un netto degrado nei modelli estremi come il 160M e il 6.9B), piuttosto questa agisce come un vincolo sul focus del modello. Infatti, osservando i modelli di taglia maggiore (da 1.4B in su), si nota che l'innesco della minaccia IT/Sistema riesce a migliorare la coerenza logit-testo rispetto alla loro baseline, forzando un maggiore allineamento tra la probabilità interna e il testo generato.

% --- 6. ALLINEAMENTO UMANO ---
\subsection{Allineamento umano}
L'\textbf{allineamento umano} quantifica il grado di somiglianza tra le opinioni espresse dal modello e la reale distribuzione delle risposte della popolazione (basata sui sondaggi del Pew Research Center). 
L'obiettivo pratico è misurare quanto le probabilità generate dalla rete si discostino dalle opinioni umane misurate empiricamente, tenendo conto della natura ordinale delle risposte (es. da "fortemente in disaccordo" a "fortemente d'accordo").

Per questa ragione, si è scelto di adottare la \textbf{Distanza di Wasserstein (WD)} \cite{santurkar2023whose}. La WD calcola il "costo" spaziale necessario per trasformare la distribuzione di probabilità del modello ($P$), basata sui suoi \textit{logit}, in quella umana ($Q$), utilizzando i pesi della scala Likert come coordinate\footnote{I pesi della scala Likert sono presi dai dati forniti in input delle risposte umane.}. 
 
Per rendere la metrica confrontabile, la distanza viene normalizzata dividendola per il valore massimo possibile ($WD_{max}$), ottenuto immaginando tutta la probabilità spostata agli estremi opposti della scala.
Lo \textit{score} di allineamento finale è definito come:
$$Score = \max\left(0, 1 - \frac{WD(P, Q)}{WD_{max}}\right)$$
I valori di questa metrica variano all'interno dell'intervallo \textbf{[0, 1]}:
\begin{itemize}
    \item \textbf{Valori prossimi allo 0}: indicano un disallineamento totale (il modello concentra la probabilità sull'estremo opposto rispetto alla popolazione);
    \item \textbf{Valori prossimi all'1}: indicano un perfetto allineamento tra la distribuzione del modello e quella umana.
\end{itemize}
Di seguito vengono riportati i risultati dell'allineamento umano per i modelli Pythia presi in considerazione, confrontando le loro distribuzioni di probabilità con quelle empiriche della popolazione.
\small
\setlength{\LTcapwidth}{\textwidth}
\begin{xltabular}{\textwidth}{X *{5}{c}}
\toprule
 & \textbf{160M} & \textbf{1B} & \textbf{1.4B} & \textbf{2.8B} & \textbf{6.9B} \\
\midrule
%% ── MACRO 6: ALLINEAMENTO UMANO ──
\rowcolor{m5hdr} \multicolumn{6}{c}{\textcolor{white}{\textbf{Allineamento umano}}} \\
\rowcolor{m5Row} \hspace{1em}Media & 0.661 & 0.787 & 0.821 & 0.809 & 0.812 \\
\rowcolor{m5Row} \hspace{1em}Mediana & 0.683 & 0.819 & 0.843 & 0.829 & 0.836 \\
\midrule
\rowcolor{m5Row} \hspace{1em}Aree più allineate & \tiny Self-p., Status & \tiny Pers.\ h., Religion & \tiny Health., Immig. & \tiny Status, Immig. & \tiny Health., Immig. \\
\rowcolor{m5Row} \hspace{1em}Aree meno allineate & \tiny Pers.\ fin., Family & \tiny Family, Job/car. & \tiny Leader., Family & \tiny Leader., Crime & \tiny Leader., Family \\
\bottomrule
\caption{Confronto dell'allineamento umano tra i cinque modelli Pythia.}
\label{tab:comparison_5models_alignment}
\end{xltabular}
\normalsize

\begin{figure}[H]
    \centering
    \includegraphics[width=\textwidth]{\risultati/comparative_paper/comp_alignment.png}
    \caption{Distribuzione dell'allineamento umano (score) per i modelli Pythia più piccoli. La linea rossa tratteggiata indica il valore medio $(\mu)$, mentre quella blu il valore mediano (med). I grafici dei modelli più grandi (2.8B e 6.9B) sono stati omessi per chiarezza visiva, in quanto presentano una distribuzione simile a quella del modello da 1.4B.}
    \label{fig:comp_alignment}
\end{figure}

I risultati sperimentali in Tabella~\ref{tab:comparison_5models_alignment} e in Figura~\ref{fig:comp_alignment} indicano che l'allineamento umano aumenta con le dimensioni del modello, ma poi si stabilizza rapidamente intorno a \textbf{1.4B} parametri, raggiungendo un valore medio di circa 0.821. Questo suggerisce che, al di là di una certa soglia dimensionale, l'aumento dei parametri non si traduce in un miglioramento significativo dell'allineamento con le opinioni umane. 

Le aree con il \textbf{maggior allineamento} sono quelle in cui le risposte umane presentano un consenso più netto: nei modelli più grandi spiccano \emph{Immigration} e \emph{Healthcare}. Al contrario, le aree con il \textbf{minor allineamento} sono \emph{Leadership}, \emph{Family} e \emph{Crime/security}, ovvero temi su cui la distribuzione delle opinioni umane è più frammentata e il modello fatica a riprodurre il pattern di risposta prevalente. L'area \emph{Leadership} è sistematicamente la più distante in tutti i modelli da 1.4B in poi, confermando che il giudizio sulla leadership politica resta uno degli aspetti più difficili da allineare per i modelli base.

% --- 7. BIAS POLITICO ---
\subsection{Affinità politica}
L'\textbf{affinità politica} è una metrica volta a classificare e a identificare quale gruppo demografico o orientamento ideologico (\textit{Democrat, Republican, Independent, Liberal, Moderate e Conservative}) sia statisticamente più vicino alle risposte generate dal modello. L'obiettivo è determinare se il sistema presenti dei bias verso una specifica fazione politica rispetto alla distribuzione reale della popolazione americana.

Anche in questo caso si usa la \textbf{Distanza di Wasserstein}. L'affinità viene calcolata individuando il gruppo $G$ che minimizza la WD tra la distribuzione delle probabilità del modello ($P$) e quella di riferimento del gruppo ($Q_G$) su un determinato set di domande:
$$\text{Affinità} = \arg\min_{G} WD(P, Q_G)$$
Trattandosi di una funzione che cerca il termine che    minimizza la distanza, il risultato non è quantitativo, ma \textbf{categoriale}, cioè restituisce direttamente l'etichetta del gruppo demografico più affine.

\textbf{Nota:} per evitare di attribuire schieramenti falsati da lievi differenze statistiche, se due o più gruppi risultano ugualmente affini in una macro-area tematica il risultato viene classificato come "Tie".

\small
\setlength{\LTcapwidth}{\textwidth}
\begin{xltabular}{\textwidth}{X *{5}{c}}
\toprule
 & \textbf{160M} & \textbf{1B} & \textbf{1.4B} & \textbf{2.8B} & \textbf{6.9B} \\
\midrule
%% ── MACRO 7: AFFINITÀ POLITICA ──
\rowcolor{m6hdr} \multicolumn{6}{c}{\textcolor{white}{\textbf{Affinità politica (\%)}}} \\
\rowcolor{m6Row} \hspace{1em}Democratici (D) & 25.4 & 20.4 & 17.3 & 17.1 & 17.3 \\
\rowcolor{m6Row} \hspace{1em}Liberali (L) & 25.0 & 14.3 & 13.0 & 13.8 & 14.8 \\
\rowcolor{m6Row} \hspace{1em}Repubblicani (R) & 25.6 & 19.0 & 19.5 & 21.5 & 18.8 \\
\rowcolor{m6Row} \hspace{1em}Conservativi (C) & 8.1 & 14.2 & 16.4 & 14.9 & 16.6 \\
\rowcolor{m6Row} \hspace{1em}Moderati (M) & 8.7 & 16.3 & 17.3 & 16.7 & 16.8 \\
\rowcolor{m6Row} \hspace{1em}Indipendenti (I) & 7.2 & 15.8 & 16.4 & 16.0 & 15.7 \\
\midrule
\rowcolor{m6Row} \hspace{1em}D + L & 50.4 & 34.7 & 30.4 & 30.8 & 32.2 \\
\rowcolor{m6Row} \hspace{1em}R + C & 33.7 & 33.2 & 35.9 & 36.5 & 35.4 \\
\bottomrule
\caption{Confronto dell'affinità politica tra i cinque modelli Pythia.}
\label{tab:comparison_5models_political_affinity}
\end{xltabular}
\normalsize
Come si nota in Tabella~\ref{tab:comparison_5models_political_affinity}, la tendenza democratica (D)/liberale (L) di Pythia-160M (D+L\,=\,50.4\,\%) si attenua progressivamente fino a 1.4B (D+L\,=\,30.4\,\%), dove avviene un \emph{cambio di rotta}: infatti, da 1.4B in poi il blocco repubblicano (R)/conservativo (C) supera quello democratico/liberale (R+C\,=\,35.9\,\%). La distribuzione si stabilizza attorno a D+L\,$\approx$\,30--32\,\% e R+C\,$\approx$\,35--36\,\%, con quote significative di moderati ($\sim$\,16\,\%) e indipendenti ($\sim$\,16\,\%).

Questa evoluzione suggerisce che i modelli diventino più \emph{neutrali} man mano che crescono di parametri. Tuttavia, questa apparente neutralità potrebbe essere interpretata in due modi: 
\begin{enumerate}
    \item Una maggiore distribuzione equa delle risposte tra i diversi orientamenti, indice di una rappresentazione più bilanciata, oppure
    \item Una maggiore \emph{confusione} o \emph{incoerenza} nelle risposte, per cui il modello non prende davvero posizione, ma distribuisce le probabilità in modo più uniforme senza una reale "comprensione" dei contenuti politici. 
\end{enumerate}
I dati sulla robustezza suggeriscono che la prima ipotesi sia più plausibile per i modelli da 1.4B in poi, dato che la JSD alla permutazione è estremamente bassa ($< 0.015$) e le risposte sono coerenti indipendentemente dall'ordine delle opzioni.

Per poter visualizzare il posizionamento ideologico dei modelli, viene utilizzata una \textbf{bussola politica}, un grafico bidimensionale che mappa ogni macro-area lungo due assi indipendenti:
\begin{enumerate}
    \item \textbf{Asse economico (sinistra/destra):} misura le opinioni sulla regolamentazione del mercato, tasse e intervento statale;
    \item \textbf{Asse sociale (autoritario/libertario):} misura le opinioni su libertà civili, diritti individuali e autorità.
\end{enumerate}
A differenza degli approcci tradizionali, le coordinate non derivano da punteggi arbitrari assegnati alle risposte (come nella scala Likert), ma dalla vicinanza statistica del modello ai gruppi umani di riferimento. 
Nello specifico, le coordinate sono ottenute calcolando il delta tra le Distanze di Wasserstein (WD) che separano le risposte del modello dai rispettivi gruppi demografici:
\begin{itemize}
    \item \textbf{Asse X (orientamento politico):} calcolato come la differenza tra la distanza dai Repubblicani e la distanza dai Democratici ($WD_{Repubblicani} - WD_{Democratici}$);
    \item \textbf{Asse Y (orientamento ideologico):} calcolato come la differenza tra la distanza dai Conservatori e la distanza dai Liberali ($WD_{Conservatori} - WD_{Liberali}$).
\end{itemize}
In questa formulazione, il grafico non misura l'accordo con specifiche tesi politiche (ad esempio libero mercato), ma mostra a quale riferimento umano la distribuzione di probabilità del modello tende ad avvicinarsi. Un valore positivo o negativo su ciascun asse segnala, quindi, una maggiore vicinanza statistica a uno dei due estremi. Il centro del grafico $(0,0)$ corrisponde a una situazione di neutralità, intesa come equidistanza dai gruppi considerati.

L'analisi si è concentrata sui due modelli che meglio rappresentano l'evoluzione del bias politico con lo \textit{scaling}.
\begin{figure}[H]
\centering
\includegraphics[width=\textwidth]{\risultati/comparative_paper/comp_political_compass.png}
\caption{Bussola politica per macro-area tematica dei modelli Pythia 160M e 1.4B: ogni punto numerato rappresenta una macro-area, posizionata in base all'affinità con i quattro gruppi politici principali.}
\label{fig:comp_political_compass}
\end{figure}

Dalla proiezione emergono due dinamiche contrapposte:
\begin{itemize}
    \item \textbf{Pythia-160M e una falsa polarizzazione:} il modello minore si posiziona in modo caotico, ma sbilanciato nel quadrante \textit{liberal/democratico}. Più che una reale preferenza ideologica, questa polarizzazione è probabile che sia un "artefatto statistico" causato dai suoi limiti strutturali (come l'estremo \textit{primacy bias} osservato in Sottosezione~\ref{subsec:permutazione}). Il modello non comprende a fondo il testo, ma seleziona meccanicamente opzioni che, nel dataset, casualmente convergono verso questo asse;
    \item \textbf{Pythia-1.4B e la tendenza alla neutralità:} con l'aumento dei parametri, il punto sulla bussola migra progressivamente verso il centro. Come evidenziato dalla stabilità nell'allineamento con le risposte umane a \textbf{0.82}, il modello 1.4B non "abbraccia" una fazione, ma impara a distribuire la probabilità in modo più simile alla popolazione generale, risultando affine a gruppi moderati e indipendenti.
\end{itemize}
L'osservazione di questa tendenza verso il centro evidenzia un risultato cruciale che si discosta dalla letteratura. Studi come quello di Santurkar et al. (2023) \cite{santurkar2023whose} hanno dimostrato che i modelli commerciali moderni possiedono forti bias \textit{progressisti} e \textit{liberal}. Tuttavia, quei sistemi sono stati sottoposti a \textit{fine-tuning} con allineamento umano (RLHF), assorbendo di fatto le preferenze e l'ideologia politica dei valutatori umani. 
I modelli Pythia analizzati in questo studio sono, al contrario, modelli \textit{pre-trained}: i risultati dimostrano che il puro \textit{scaling} architetturale sui dati non pre-filtrati di internet non genera derive ideologiche verso sinistra, ma tende anzi a stabilizzare il sistema su una rappresentazione variata e bilanciata delle opinioni umane, evitando delle polarizzazioni nette.

\end{document}